% Kooperationsvereinbarung für ein Pilotprojekt mit Anspruch.
% ENTWURF ZUR ANWALTLICHEN PRÜFUNG. Keine Rechtsberatung.
% Quellen: docs/proposal-organisations.md (Angebot, Bitte, Kosten, Lizenz,
% Grenzen), docs/PARTNERS.md (Partnerlink, Summen, Übersicht oder monatliche
% Zusammenstellung), docs/COMPLIANCE.md, lib/legal/datenschutz.ts
% (Datenschutzerklärung), README.md (Grenzen, Lizenz offen).
% Genannte Normen: nur § 7a SGB XI. Weitere Normen setzt die anwaltliche
% Prüfung ein; jede solche Stelle trägt ein \todo{}.
% Lokale Makros: \fillin (Kasten mit den Platzhaltern), \sigfield
% (Unterschriftsfeld), Kopfzeile mit "Entwurf zur anwaltlichen Prüfung"
% (überschreibt \ihead und \ohead der Präambel nur in diesem Dokument).
\documentclass[fontsize=11pt,paper=a4,parskip=half]{scrartcl}
% Shared preamble for the Anspruch dossier and its companion documents.
% Load after \documentclass with \input{preamble} (from docs/dossier/)
% or % Shared preamble for the Anspruch dossier and its companion documents.
% Load after \documentclass with \input{preamble} (from docs/dossier/)
% or % Shared preamble for the Anspruch dossier and its companion documents.
% Load after \documentclass with \input{preamble} (from docs/dossier/)
% or \input{../preamble} (from docs/dossier/antrag/ or docs/dossier/partner/).
% Each document sets its own \documentclass and its own pdftitle via
% \hypersetup{pdftitle={...}}. Only packages that ship with MiKTeX.

\usepackage[T1]{fontenc}
\usepackage[utf8]{inputenc}
\usepackage{lmodern}
\usepackage[english,ngerman]{babel} % last option = main language: ngerman
\usepackage[a4paper,margin=25mm]{geometry}
\usepackage{booktabs}
\usepackage{longtable}
\usepackage{array}
\usepackage{enumitem}
\usepackage{xcolor}
\usepackage{microtype}
\usepackage[autostyle=true,german=quotes]{csquotes}
\usepackage{marginnote}
\usepackage[hidelinks,unicode,pdflang={de-DE}]{hyperref}

% ---------------------------------------------------------------------------
% Colours
\definecolor{blankbg}{RGB}{255,236,153}   % highlighted placeholder
\definecolor{todofg}{RGB}{170,40,40}      % "to confirm" margin note

% \blank{Text}: a highlighted placeholder "[Text]" for values the sender
% fills in before sending. Usable in running text and in table cells.
\newcommand{\blank}[1]{\begingroup\setlength{\fboxsep}{1pt}%
  \colorbox{blankbg}{\strut[#1]}\endgroup}

% \todo{Text}: a "zu bestätigen" note in the margin. Works in running text,
% lists and longtable cells (via marginnote, not \marginpar).
% Defined with \providecommand so that it does not clash if a document
% loads todonotes on its own (todonotes then wins).
\providecommand{\todo}[1]{\marginnote{\raggedright\scriptsize
  \textcolor{todofg}{\textbf{Zu bestätigen:} #1}}}
\setlength{\marginparwidth}{20mm}

% \draftmark: header text for every page of a draft.
\newcommand{\draftmark}{Entwurf, 23.10.2026}

% Draft header for scrartcl / scrreprt (KOMA) and article / report.
% Letters (scrlttr2, letter) place \draftmark themselves.
\makeatletter
\@ifclassloaded{scrartcl}{\def\anspruch@koma{}}{}
\@ifclassloaded{scrreprt}{\def\anspruch@koma{}}{}
\@ifclassloaded{article}{\def\anspruch@std{}}{}
\@ifclassloaded{report}{\def\anspruch@std{}}{}
\ifdefined\anspruch@koma
  \usepackage[headsepline]{scrlayer-scrpage}
  \clearpairofpagestyles
  \ohead{\small\draftmark}
  \ihead{\small\headmark}
  \cfoot*{\pagemark}
  \automark[section]{section}
  \pagestyle{scrheadings}
\fi
\ifdefined\anspruch@std
  \usepackage{fancyhdr}
  \pagestyle{fancy}
  \fancyhf{}
  \fancyhead[R]{\small\draftmark}
  \fancyhead[L]{\small\nouppercase{\leftmark}}
  \fancyfoot[C]{\thepage}
  \fancypagestyle{plain}{\fancyhf{}\fancyhead[R]{\small\draftmark}%
    \fancyfoot[C]{\thepage}\renewcommand{\headrulewidth}{0pt}}
\fi
\makeatother

% Lists: compact, as in a German Antrag.
\setlist{itemsep=0.2em,topsep=0.4em}

% Tables: a bit more air, ragged-right paragraph columns.
\renewcommand{\arraystretch}{1.15}
\newcolumntype{L}[1]{>{\raggedright\arraybackslash}p{#1}}
 (from docs/dossier/antrag/ or docs/dossier/partner/).
% Each document sets its own \documentclass and its own pdftitle via
% \hypersetup{pdftitle={...}}. Only packages that ship with MiKTeX.

\usepackage[T1]{fontenc}
\usepackage[utf8]{inputenc}
\usepackage{lmodern}
\usepackage[english,ngerman]{babel} % last option = main language: ngerman
\usepackage[a4paper,margin=25mm]{geometry}
\usepackage{booktabs}
\usepackage{longtable}
\usepackage{array}
\usepackage{enumitem}
\usepackage{xcolor}
\usepackage{microtype}
\usepackage[autostyle=true,german=quotes]{csquotes}
\usepackage{marginnote}
\usepackage[hidelinks,unicode,pdflang={de-DE}]{hyperref}

% ---------------------------------------------------------------------------
% Colours
\definecolor{blankbg}{RGB}{255,236,153}   % highlighted placeholder
\definecolor{todofg}{RGB}{170,40,40}      % "to confirm" margin note

% \blank{Text}: a highlighted placeholder "[Text]" for values the sender
% fills in before sending. Usable in running text and in table cells.
\newcommand{\blank}[1]{\begingroup\setlength{\fboxsep}{1pt}%
  \colorbox{blankbg}{\strut[#1]}\endgroup}

% \todo{Text}: a "zu bestätigen" note in the margin. Works in running text,
% lists and longtable cells (via marginnote, not \marginpar).
% Defined with \providecommand so that it does not clash if a document
% loads todonotes on its own (todonotes then wins).
\providecommand{\todo}[1]{\marginnote{\raggedright\scriptsize
  \textcolor{todofg}{\textbf{Zu bestätigen:} #1}}}
\setlength{\marginparwidth}{20mm}

% \draftmark: header text for every page of a draft.
\newcommand{\draftmark}{Entwurf, 23.10.2026}

% Draft header for scrartcl / scrreprt (KOMA) and article / report.
% Letters (scrlttr2, letter) place \draftmark themselves.
\makeatletter
\@ifclassloaded{scrartcl}{\def\anspruch@koma{}}{}
\@ifclassloaded{scrreprt}{\def\anspruch@koma{}}{}
\@ifclassloaded{article}{\def\anspruch@std{}}{}
\@ifclassloaded{report}{\def\anspruch@std{}}{}
\ifdefined\anspruch@koma
  \usepackage[headsepline]{scrlayer-scrpage}
  \clearpairofpagestyles
  \ohead{\small\draftmark}
  \ihead{\small\headmark}
  \cfoot*{\pagemark}
  \automark[section]{section}
  \pagestyle{scrheadings}
\fi
\ifdefined\anspruch@std
  \usepackage{fancyhdr}
  \pagestyle{fancy}
  \fancyhf{}
  \fancyhead[R]{\small\draftmark}
  \fancyhead[L]{\small\nouppercase{\leftmark}}
  \fancyfoot[C]{\thepage}
  \fancypagestyle{plain}{\fancyhf{}\fancyhead[R]{\small\draftmark}%
    \fancyfoot[C]{\thepage}\renewcommand{\headrulewidth}{0pt}}
\fi
\makeatother

% Lists: compact, as in a German Antrag.
\setlist{itemsep=0.2em,topsep=0.4em}

% Tables: a bit more air, ragged-right paragraph columns.
\renewcommand{\arraystretch}{1.15}
\newcolumntype{L}[1]{>{\raggedright\arraybackslash}p{#1}}
 (from docs/dossier/antrag/ or docs/dossier/partner/).
% Each document sets its own \documentclass and its own pdftitle via
% \hypersetup{pdftitle={...}}. Only packages that ship with MiKTeX.

\usepackage[T1]{fontenc}
\usepackage[utf8]{inputenc}
\usepackage{lmodern}
\usepackage[english,ngerman]{babel} % last option = main language: ngerman
\usepackage[a4paper,margin=25mm]{geometry}
\usepackage{booktabs}
\usepackage{longtable}
\usepackage{array}
\usepackage{enumitem}
\usepackage{xcolor}
\usepackage{microtype}
\usepackage[autostyle=true,german=quotes]{csquotes}
\usepackage{marginnote}
\usepackage[hidelinks,unicode,pdflang={de-DE}]{hyperref}

% ---------------------------------------------------------------------------
% Colours
\definecolor{blankbg}{RGB}{255,236,153}   % highlighted placeholder
\definecolor{todofg}{RGB}{170,40,40}      % "to confirm" margin note

% \blank{Text}: a highlighted placeholder "[Text]" for values the sender
% fills in before sending. Usable in running text and in table cells.
\newcommand{\blank}[1]{\begingroup\setlength{\fboxsep}{1pt}%
  \colorbox{blankbg}{\strut[#1]}\endgroup}

% \todo{Text}: a "zu bestätigen" note in the margin. Works in running text,
% lists and longtable cells (via marginnote, not \marginpar).
% Defined with \providecommand so that it does not clash if a document
% loads todonotes on its own (todonotes then wins).
\providecommand{\todo}[1]{\marginnote{\raggedright\scriptsize
  \textcolor{todofg}{\textbf{Zu bestätigen:} #1}}}
\setlength{\marginparwidth}{20mm}

% \draftmark: header text for every page of a draft.
\newcommand{\draftmark}{Entwurf, 23.10.2026}

% Draft header for scrartcl / scrreprt (KOMA) and article / report.
% Letters (scrlttr2, letter) place \draftmark themselves.
\makeatletter
\@ifclassloaded{scrartcl}{\def\anspruch@koma{}}{}
\@ifclassloaded{scrreprt}{\def\anspruch@koma{}}{}
\@ifclassloaded{article}{\def\anspruch@std{}}{}
\@ifclassloaded{report}{\def\anspruch@std{}}{}
\ifdefined\anspruch@koma
  \usepackage[headsepline]{scrlayer-scrpage}
  \clearpairofpagestyles
  \ohead{\small\draftmark}
  \ihead{\small\headmark}
  \cfoot*{\pagemark}
  \automark[section]{section}
  \pagestyle{scrheadings}
\fi
\ifdefined\anspruch@std
  \usepackage{fancyhdr}
  \pagestyle{fancy}
  \fancyhf{}
  \fancyhead[R]{\small\draftmark}
  \fancyhead[L]{\small\nouppercase{\leftmark}}
  \fancyfoot[C]{\thepage}
  \fancypagestyle{plain}{\fancyhf{}\fancyhead[R]{\small\draftmark}%
    \fancyfoot[C]{\thepage}\renewcommand{\headrulewidth}{0pt}}
\fi
\makeatother

% Lists: compact, as in a German Antrag.
\setlist{itemsep=0.2em,topsep=0.4em}

% Tables: a bit more air, ragged-right paragraph columns.
\renewcommand{\arraystretch}{1.15}
\newcolumntype{L}[1]{>{\raggedright\arraybackslash}p{#1}}

\hypersetup{pdftitle={Kooperationsvereinbarung Pilotprojekt Anspruch (Entwurf)}}

% Kopfzeile: sichtbarer Entwurfsvermerk auf jeder Seite.
\ihead{\small\textbf{Entwurf zur anwaltlichen Prüfung}}
\ohead{\small\draftmark}
% Paragrafen statt Abschnittsnummern.
\renewcommand*{\thesection}{§\,\arabic{section}}
\renewcommand*{\sectionformat}{\thesection\enskip}
\RedeclareSectionCommand[beforeskip=1.2ex plus .2ex,afterskip=.4ex]{section}
\setkomafont{section}{\normalfont\large\bfseries}

% parskip=half reserves 1em at the end of a paragraph; the full-width box needs
% that switched off locally.
\newcommand{\fillin}[2]{{\setlength{\parfillskip}{0pt plus 1fil}\noindent
  \fbox{\parbox{\dimexpr\linewidth-2\fboxsep-2\fboxrule\relax}{%
  \footnotesize\raggedright\textbf{#1} #2}}\par}}
\newcommand{\sigfield}[1]{\parbox[t]{0.45\linewidth}{\vspace{14mm}\hrule\vspace{1mm}%
  \footnotesize #1}}

\begin{document}

\fillin{Vor dem Versand ausfüllen:}{Name des Betreibers; Anschrift des
Betreibers; Organisation; Anschrift der Organisation; vertreten durch (Name,
Funktion); Startdatum; Adresse der Seite; kürzel; Betrag; Zahlungsweise;
Kündigungsfrist. Nach der anwaltlichen Prüfung: Regelung zur Haftung; anwendbares
Recht und Gerichtsstand; alle Stellen mit \enquote{Zu bestätigen} am Rand.}

\begin{center}
  {\Large\bfseries Kooperationsvereinbarung}\\[0.5ex]
  {\large über ein Pilotprojekt mit der Webseite \enquote{Anspruch}}\\[0.5ex]
  {\small\textbf{Entwurf zur anwaltlichen Prüfung.} Noch nicht zur Unterschrift.}
\end{center}

\section{Parteien}
Diese Vereinbarung schließen
\begin{enumerate}[label=(\arabic*)]
  \item \blank{Name des Betreibers}, \blank{Anschrift des Betreibers},
    Betreiber der Webseite Anspruch (im Folgenden \enquote{Betreiber}), und
  \item \blank{Organisation}, \blank{Anschrift der Organisation}, vertreten durch
    \blank{Name}, \blank{Funktion} (im Folgenden \enquote{Partner}).
\end{enumerate}

\section{Gegenstand}
(1) Anspruch ist eine Webseite für Familien, die einen Angehörigen pflegen. Die
Familie beantwortet Fragen zum Alltag der pflegebedürftigen Person. Die Seite
schätzt den Pflegegrad, den die Begutachtung wahrscheinlich ergibt. Sie vergleicht
die Leistungen, die daraus folgen, mit dem, was der Haushalt schon bekommt, und
nennt die nächsten Schritte mit der gesetzlichen Grundlage.

(2) Gegenstand dieser Vereinbarung ist ein Pilotprojekt. Der Partner erprobt die
Seite in seiner Beratung. Die Parteien werten die Erfahrungen gemeinsam aus.

\section{Laufzeit}
(1) Das Pilotprojekt beginnt am \blank{Startdatum} und dauert drei Monate.

(2) Es endet nach drei Monaten, ohne dass es einer Kündigung bedarf. Eine
Verlängerung bedarf einer schriftlichen Vereinbarung.

\section{Leistungen des Betreibers}
Der Betreiber
\begin{enumerate}[label=(\arabic*)]
  \item stellt die Webseite unter \blank{Adresse der Seite} für die Laufzeit
    bereit. Familien brauchen kein Konto.\todo{Ob eine Verfügbarkeit zugesagt
    wird und in welchem Umfang.}
  \item richtet für den Partner einen eigenen Link ein:
    \texttt{https://}\blank{Adresse der Seite}\texttt{/?via=}\blank{kürzel}.
    Besuche über diesen Link werden der Organisation zugeordnet, nicht einer
    Person.
  \item stellt dem Partner die Zahlen zu diesem Link zur Verfügung, nur als
    Summen. Das geschieht per Link zu einer gefilterten Übersicht oder als
    monatliche Zusammenstellung. Die Zahlen sind: Besuche; wie viele begonnen
    haben; wie viele ein Ergebnis erreicht haben; wie viele ein PDF geladen
    haben; die Sprache der Oberfläche; die Spanne der gestellten Fragen; die
    Spanne, in der die geschätzte monatliche Lücke lag. Voraussetzung ist, dass
    die Zählung eingeschaltet ist.\todo{Ob die Zählung zum Start eingeschaltet
    ist und welcher der beiden Wege gilt.}
  \item führt einmal im Monat ein Rückmeldegespräch mit dem Partner.
  \item stellt dem Partner das Factsheet und den Vorschlag für eine
    Pilotpartnerschaft zur Verfügung.
  \item teilt dem Partner vor Beginn des Pilotprojekts mit, unter welcher Lizenz
    der Quellcode steht (siehe §\,10).
\end{enumerate}

\section{Leistungen des Partners}
Der Partner
\begin{enumerate}[label=(\arabic*)]
  \item benennt zwei Beraterinnen oder Berater. Sie nutzen die Seite im Gespräch
    oder geben Familien den Link mit.
  \item nimmt einmal im Monat am Rückmeldegespräch teil.
  \item lässt seine Beratenden eine Strichliste ohne Namen führen: wie oft die
    Seite genutzt wurde und wie oft danach ein Antrag, ein Antrag auf
    Höherstufung oder ein Widerspruch gestellt wurde. Die Strichliste enthält
    keine Angaben, die sich einer Person zuordnen lassen.
  \item gibt den Link nur im eigenen Namen weiter.
\end{enumerate}

\section{Datenschutz}
(1) Die Parteien tauschen keine personenbezogenen Daten aus. Das gilt auch für die
Strichliste nach §\,5 Abs.\,3.

(2) Was eine Familie auf der Seite tippt, bleibt auf ihrem Gerät. Die Antworten
werden nur in ihrem Browser gespeichert und lassen sich dort mit einem Tipp löschen.
Der Betreiber erhält sie nicht. Antwortet die Familie per Sprache, erkennt der
Browser die Sprache, nicht die Seite. Dafür gilt die Datenschutzerklärung des
Browserherstellers.

(3) Ist die Zählung eingeschaltet, werden nur Besuche und fünf Schritte gezählt
(Start, Ergebnis, weitere Fragen, PDF, Löschen). Die Zählung arbeitet ohne Cookies,
ohne Namen, ohne Antworten und ohne Pflegegrad. Das Kürzel des Links bezeichnet die
Organisation, nicht eine Person. Der Partner erhält nur Summen.

(4) Jede Partei ist für ihre eigene Verarbeitung personenbezogener Daten selbst
verantwortlich. Der Partner für seine Beratung, der Betreiber für die Webseite nach
ihrer Datenschutzerklärung.

(5) Der Betreiber verarbeitet keine personenbezogenen Daten im Auftrag des Partners.
Eine Auftragsverarbeitung entsteht nicht.\todo{Rollen nach der DSGVO (Abs.~4
und 5): keine gemeinsame Verantwortung, keine Auftragsverarbeitung, kein
Vertrag dazu nötig.}

\section{Charakter des Werkzeugs}
(1) Die Seite liefert eine Schätzung aus den Angaben der Familie. Sie ist keine
Begutachtung durch den Medizinischen Dienst und keine Rechtsberatung.

(2) Die Seite sagt jeder Familie, dass ihr die kostenlose Pflegeberatung nach
§\,7a SGB~XI zusteht. Sie ersetzt diese Beratung nicht. Die Beratenden des Partners
bleiben für ihre Beratung selbst verantwortlich.

(3) Die Beträge und Punkteschwellen folgen dem geltenden Recht. Sie wurden zuletzt
am 22.10.2026 geprüft. Das Pflegeneuordnungsgesetz ist ein Kabinettsentwurf vom
30.09.2026 und noch kein Gesetz. Wird es wie entworfen beschlossen, ändern sich die
meisten Beträge ab dem 01.01.2027 und die Seite muss angepasst werden.

(4) Die türkische, russische, polnische und arabische Oberfläche ist noch nicht von
Muttersprachlern geprüft, sondern bisher nur von einem Sprachmodell gegen den
deutschen Text.

\section{Haftung}
\blank{Regelung zur Haftung}\todo{Haftung beider Parteien regeln, insbesondere
für eine unzutreffende Schätzung und für Ausfälle der Seite.}

\section{Kosten}
(1) Für das Pilotprojekt zahlt der Partner \blank{Betrag}.

(2) Zahlungsweise: \blank{Zahlungsweise}.

\section{Nutzungsrechte}
(1) Der Partner darf den Link nach §\,4 Abs.\,2 und die Unterlagen nach §\,4 Abs.\,5
für die Dauer des Pilotprojekts in seiner Beratung und in Hinweisen an Familien
nutzen.

(2) Der Quellcode der Seite hat noch keine Lizenz. Der Betreiber teilt die
Entscheidung vor Beginn des Pilotprojekts mit. Diese Vereinbarung räumt keine
Rechte am Quellcode ein.\todo{Lizenz des Quellcodes und ihre Folgen für den
Partner.}

(3) Der Betreiber nennt den Partner nur mit dessen schriftlicher Zustimmung als
Pilotpartner.

\section{Kündigung}
(1) Jede Partei kann diese Vereinbarung mit einer Frist von \blank{Kündigungsfrist}
schriftlich kündigen.

(2) Das Recht zur Kündigung aus wichtigem Grund bleibt unberührt.\todo{Folgen
einer Kündigung für bereits gezahlte Kosten.}

\section{Schlussbestimmungen}
(1) Änderungen und Ergänzungen bedürfen der Schriftform.

(2) Ist eine Bestimmung unwirksam, bleibt der Rest der Vereinbarung wirksam.

(3) Anwendbares Recht und Gerichtsstand: \blank{anwendbares Recht und
Gerichtsstand}.\todo{Rechtswahl und Gerichtsstand festlegen; Fassung von
Abs.~2.}

\vspace{6mm}
\noindent\sigfield{Ort, Datum\\Unterschrift\\\blank{Name des Betreibers}, Betreiber}\hfill
\sigfield{Ort, Datum\\Unterschrift und Stempel\\\blank{Name}, \blank{Funktion},
\blank{Organisation}}

\end{document}
